\section{Additional Exercises}

\subsection{Recognising \& Dealing With Different Signatures}

\begin{exercise}\label{ex:metricmfds-1}
(reworded) State the possible signatures of a metric $g$ for the following set of $g$-null vectors:
\benr
    \item A cone through the origin,
    \item A point at the origin,
    \item A straight line through the origin, and
    \item A plane through the origin,
\een
where the origin is the origin of the vector space, namely the point $p\in\cM$ that we are tangent to.
\end{exercise}
\hyperref[sol:metricmfds-1]{\small \textit{Solution on p.~\pageref*{sol:metricmfds-1}}}

\subsection{Levi-Civita Connection}

\begin{exercise}\label{ex:metricmfds-2}
Expand in terms of connection coefficient functions
\benr
    \item $\big(\nabla_ag\big)_{bc}$,
    \item $\big(\nabla_bg\big)_{ca}$,
    \item $\big(\nabla_cg\big)_{ab}$.
\een
\end{exercise}
\hyperref[sol:metricmfds-2]{\small \textit{Solution on p.~\pageref*{sol:metricmfds-2}}}

\begin{exercise}\label{ex:metricmfds-3}
By adding and/or subtracting (i), (ii) and (iii) above in a clever way, obtain
\bse
    {\Gamma^a}_{bc} = \frac{1}{2}\big(g^{-1}\big)^{am} \big( g_{mc,b} + g_{mb,c} - g_{bc,m}\big)
\ese
and conclude that $\nabla g=0$ and $T=0$ (torsion) uniquely determine the connection coefficient functions in terms of the metric.
\end{exercise}
\hyperref[sol:metricmfds-3]{\small \textit{Solution on p.~\pageref*{sol:metricmfds-3}}}

\subsection{Massaging The Length Functional}

\begin{exercise}\label{ex:metricmfds-4}
Let $\gamma :(0,1) \to \cM$ be a smooth curve on a smooth manifold $(\cM,\cO,\cA)$. Now consider a second curve $\widetilde{\gamma}:(0,1)\to\cM$ defined by
\bse
    \widetilde{\gamma}(\lambda) = \gamma\big(\sig(\lambda)\big),
\ese
where $\sig:(0,1)\to(0,1)$ is an increasing bijective smooth function.

Show that the length of both curves is the same:
\bse
    L[\widetilde{\gamma}] = L[\gamma].
\ese
\end{exercise}
\hyperref[sol:metricmfds-4]{\small \textit{Solution on p.~\pageref*{sol:metricmfds-4}}}

\begin{exercise}\label{ex:metricmfds-5}
Show that the Euler-Lagrange equations for a Lagrangian $\cT$ have precisely the same solutions as the Euler-Lagrange equations for the Lagrangian $\cL :=\sqrt{\cT}$, if of the latter one only selects those solutions that satisfy the condition $\cT=1$ on their parameterisation.
\end{exercise}
\hyperref[sol:metricmfds-5]{\small \textit{Solution on p.~\pageref*{sol:metricmfds-5}}}

\subsection{A Practical Way To Quickly Determine Christoffel Symbols}

\begin{exercise}\label{ex:metricmfds-6}
Derive the geodesic equation for the two-dimensional round sphere of radius $R$, whose metric in some chart $(U,x)$ is given by
\bse
    g_{ab}\big( x^{-1}(\vartheta,\phi)\big) = \begin{pmatrix}
        R^2 & 0 \\
        0 & R^2 \sin^2\vartheta
    \end{pmatrix}
\ese
via a convenient Euler-Lagrange equation. In order to lighten notation, you may define
\bse
    \vartheta(\lambda) := (x^1\circ\gamma)(\lambda), \qand \phi(\lambda) := (x^2\circ\gamma)(\lambda).
\ese
\end{exercise}
\hyperref[sol:metricmfds-6]{\small \textit{Solution on p.~\pageref*{sol:metricmfds-6}}}

\begin{exercise}\label{ex:metricmfds-7}
Read off the metric-induced connection coefficient functions for the round sphere.
\end{exercise}
\hyperref[sol:metricmfds-7]{\small \textit{Solution on p.~\pageref*{sol:metricmfds-7}}}

\subsection{Properties Of The Riemann-Christoffel Tensor}

\begin{exercise}\label{ex:metricmfds-8}
Show that the chart-induced basis fields act on the coefficient functions as
\bse
    \frac{\p}{\p x^c}\big(g^{-1}\big)^{ab} = - \big(g^{-1}\big)^{ar} \big(g^{-1}\big)^{bs} \frac{\p}{\p x^c}g_{rs}.
\ese
\end{exercise}
\hyperref[sol:metricmfds-8]{\small \textit{Solution on p.~\pageref*{sol:metricmfds-8}}}

\begin{exercise}\label{ex:metricmfds-9}
Use normal coordinates to find an expression for the Riemann-Christoffel tensor
\bse
    R_{abcd} = g_{ak}{R^k}_{bcd}
\ese
at a given point $p$ in terms of $g_{ab}$ and its first and second order derivatives at that very point.
\end{exercise}
\hyperref[sol:metricmfds-9]{\small \textit{Solution on p.~\pageref*{sol:metricmfds-9}}}

\begin{exercise}\label{ex:metricmfds-10}
Show --- in normal coordinates --- that $R_{abcd}=-R_{bacd}$.
\end{exercise}
\hyperref[sol:metricmfds-10]{\small \textit{Solution on p.~\pageref*{sol:metricmfds-10}}}

\begin{exercise}\label{ex:metricmfds-11}
Similarly, show that $R_{abcd}=R_{cdab}$.
\end{exercise}
\hyperref[sol:metricmfds-11]{\small \textit{Solution on p.~\pageref*{sol:metricmfds-11}}}

\begin{exercise}\label{ex:metricmfds-12}
Show that $R_{a[bcd]}=0$ for the Riemann-Christoffel tensor.
\end{exercise}
\hyperref[sol:metricmfds-12]{\small \textit{Solution on p.~\pageref*{sol:metricmfds-12}}}
